\section{Padding from a bounded degree preserves full multiplicity}

\subsection{O-012.a: a determinant bound obtained by splitting degree}

For $b\geq1$ and $r\geq2$, set
\[
 u=\left\lfloor\frac{b-1}{2}\right\rfloor,\qquad
 v=\left\lceil\frac{b-1}{2}\right\rceil,\qquad
 B(b,r)=\binom{r-1+u}{u}+\binom{r-1+v}{v}-1.
\]
Every degree-$b$ homogeneous form $f$ in $r$ variables, with a
chosen nonzero linear form $L$, admits a homogeneous linear
determinant of size $B(b,r)$ equal to
\[
                           L^{B(b,r)-b}f.               \tag{O-012.1}
\]
Consequently $L^{n-b}f\in\Omega_n$ whenever $n\geq B(b,r)$
and $r\leq n^2$. In one variable a diagonal determinant of size
$b$ suffices. These are integral identities in universal
coefficient variables.

\begin{proof}
Use $L$ and $z_1,\ldots,z_{r-1}$ as coordinates, and write
$f=\sum_{|\gamma|\leq b}c_\gamma z^\gamma L^{b-|\gamma|}$.
Make two copies of the monomial tree on the \emph{same} $r-1$
variables. The first has vertices of total degree at most $u$
and edges pointing away from its root; the second has vertices
of total degree at most $v$ and edges pointing toward its root.
The parent of a nonzero multi-index is obtained by subtracting
the last nonzero coordinate, so paths are unique. Label the
variable edges by $-z_j$, and identify the two roots into one
central vertex. Keep all other vertices distinct, even when
their multi-index labels are the same. There are $B=B(b,r)$ vertices.

Fix once and for all a rule that splits every multi-index of
degree at most $u+v=b-1$ as $\alpha+\beta$, with
$|\alpha|\leq u$, $|\beta|\leq v$; taking an initial segment
of its sorted variable word gives such a rule. For each term
with $|\gamma|\leq b-1$, insert the coefficient edge from
the prefix vertex $\alpha$ to the suffix vertex $\beta$ with
label $c_\gamma L$. For $|\gamma|=b$, remove one variable
$z_j$ by a fixed rule, split $\gamma-e_j=\alpha+\beta$,
and put label $c_\gamma z_j$ on that coefficient edge.
If several labels have the same endpoints, add them in the
matrix entry. All entries remain homogeneous linear forms.
Put an $L$-loop at every noncentral vertex; the central entry
is exactly the sum of its coefficient-loop labels.

As in M-007, every non-loop directed cycle uses the central
vertex and exactly one coefficient edge. Thus a determinant
term chooses one coefficient-edge summand and the unique
variable paths to and from it, with loops at all unused vertices.
For $|\gamma|=k<b$, the cycle has $k+1$ edges. Its sign
$(-1)^k$ cancels the sign on the $k$ variable edges, and its
contribution is $c_\gamma z^\gamma L^{B-k}$.
For $|\gamma|=b$, the cycle has $b$ edges, of which $b-1$
are tree edges and the last carries $c_\gamma z_j$.
The signs again cancel and the contribution is
$c_\gamma z^\gamma L^{B-b}$.
Summation gives (O-012.1). Since $r\geq2$, the two trees have
at least $u+1$ and $v+1$ vertices, so $B\geq u+v+1=b$;
all exponents are nonnegative. Adjoin the block $L I_{n-B}$
when a larger size $n$ is required.
\end{proof}

This complements the variable-splitting bound M-007. Splitting
degree is more efficient when the base degree is fixed and the
number of variables grows. One may use the smaller of the two
proved bounds. No assertion that either is an optimal general
determinantal representation is made.

\subsection{O-012.b: an arbitrary-rank padding inequality}

Let $U$ have dimension $r$, let $b\leq n$, and let
$\mu\vdash db$ have length at most $r$. Put
\[
                \lambda=\mu+d(n-b)e_1.
\]
If a closed $\mathrm{GL}(V)$-stable cone
$X\subseteq\operatorname{Sym}^n V^*$ contains every form
$x_1^{n-b}f$ with $f\in\operatorname{Sym}^b U^*$, then
\[
       \operatorname{mult}_{S_\lambda V}\mathbb C[X]_d
                                     \geq a_\mu(d[b]).  \tag{O-012.2}
\]
Here $U$ is the span of the first $r$ coordinate directions
in $V$. This is an arbitrary-rank form of the padding inequality
stated in Ikenmeyer--Panova Proposition~2.6(b), p.~6, citing
Kadish--Landsberg Theorem~1.3. The proof below makes the
normalization issue explicit.

\begin{proof}
Put $p=n-b$, and let
$T:\operatorname{Sym}^b U\to\operatorname{Sym}^n U$
be multiplication by $e_1^p$. For the upper Borel,
$gT=g_{11}^pTg$, so $\operatorname{Sym}^d T$ takes a
highest-weight vector of weight $\mu$ to one of weight
$\lambda$. It is injective.

Use the differential pairings
$\langle a,F\rangle=a(\partial)F|_0$ in each symmetric power.
The transpose $R$ of the padding map
$P:f\mapsto x_1^pf$ is $\partial_{e_1}^p$. On a monomial
$e^\alpha$ of degree $b$ one has
\[
                    RT(e^\alpha)
                    =\frac{(\alpha_1+p)!}{\alpha_1!}e^\alpha.
\]
Thus $RT$ is an invertible diagonal operator in characteristic
zero. Pullback to the padded linear space, after applying
$\operatorname{Sym}^dT$, is
$\operatorname{Sym}^d(RT)$, also invertible.
In particular, a basis of the degree-$d$ highest-weight space
of weight $\mu$ lifts to highest-weight vectors of weight
$\lambda$ whose restrictions to that padded space remain
linearly independent. Since that space lies in $X$, the lifted
vectors remain independent in $\mathbb C[X]_d$, proving
(O-012.2). A different conventional normalization of the
symmetric-power pairing changes the displayed diagonal factors
by nonzero constants and does not change this argument.
\end{proof}

Together O-012.a and O-012.b give
\[
             K_n(\mu+d(n-b)e_1)\geq a_\mu(d[b])
                 \quad\text{if }n\geq B(b,r).           \tag{O-012.3}
\]
The geometric statement concerns a whole padded linear space;
it does not assert identical evaluation of the lifted polynomial
on a single unchanged base form. Invertibility of the diagonal
pullback, rather than a false coefficient-normalization identity,
is what preserves all copies.

\subsection{O-012.c: stabilization after removing the last-row rectangle}

Let $d\geq2$ and $\lambda\vdash nd$ have at most $d$ rows;
pad it with zeros to length $d$. Define
\[
 k=\lambda_d,\qquad
 \rho=(\lambda_2-k,\ldots,\lambda_d-k),\qquad
 t=|\rho|,\qquad b=k+t.
\]
Zero parts of $\rho$ may be deleted. If $b\geq1$ and $n\geq b$,
put
\[
       \mu=(k+(d-1)t,\lambda_2,\ldots,\lambda_d)\vdash db.
\]
Then
\[
                         a_\lambda(d[n])=a_\mu(d[b]).  \tag{O-012.4}
\]
When $b=0$, the partition is the one-row shape and its ambient
multiplicity is one. For $\ell(\lambda)<d$ the definition gives
$k=0$, $b=|\bar\lambda|$; for $\ell(\lambda)=d$ it can
give a much smaller base degree than the total tail size.

\begin{proof}
Compute stable plethysm multiplicities using a $d$-dimensional
row space and a $d$-dimensional column space, with matrix
variables $Z=(z_{ij})$. By the Cauchy decomposition, the row
highest-weight space of weight $\lambda$ is the column module
$S_\lambda\mathbb C^d$. Its subspace of column multidegree
$(n,\ldots,n)$ fixed by column permutations has dimension
$a_\lambda(d[n])$.

Every polynomial in this row highest-weight space is divisible
by $\det(Z)^k$. To justify the assertion about every vector,
multiplication by that determinant power embeds the row
highest-weight space of weight $\lambda-k^d$ into the one
of weight $\lambda$. As column modules the map is the nonzero
map between the identical irreducibles
$\det^k\otimes S_{\lambda-k^d}\mathbb C^d$ and
$S_\lambda\mathbb C^d$, so it is an isomorphism.

After division, write $a=n-k$. The row weight is
$(da-t,\rho)$ and every column has degree $a$.
Each monomial has total degree $t$ outside the first row.
Consequently its exponent of $z_{1j}$ is at least $a-t$
in every column $j$. Since $a\geq t$, every such polynomial
is divisible by
\[
                         Q=\prod_{j=1}^d z_{1j}^{a-t}.
\]
This factor is fixed by the upper row unipotent group and by
column permutations. Division by it preserves the highest-weight
condition, subtracts $d(a-t)e_1$ from the row weight, and
reduces every column degree to $t$. Conversely multiplication
by $Q$ reverses the operation. Thus it is a bijection of the
two relevant column-permutation modules, including all their
copies and their signs.

Multiplying back by $\det(Z)^k$ gives row weight
$(k+(d-1)t,\lambda_2,\ldots,\lambda_d)=\mu$ and
column degree $b=k+t$. The determinant factor has the same
column sign character before and after, so the invariant
dimensions agree. This proves (O-012.4). The tuple $\mu$ is
a partition: $\rho_1\leq t\leq(d-1)t$, so its first
row is at least its second. Cauchy naturality then carries
the stable multiplicity to every actual vector-space dimension
in which $\lambda$ and $\mu$ fit.
\end{proof}

\subsection{O-012.d: full actual multiplicities and the four-copy family}

Under the hypotheses of O-012.c, set $r=\ell(\lambda)$.
If $r\leq n^2$ and
\[
                          n\geq\max\{b,B(b,r)\},       \tag{O-012.5}
\]
then
\[
                         K_n(\lambda)=a_\lambda(d[n]). \tag{O-012.6}
\]
For $r=1$ use the pure-power diagonal determinant instead of
$B(b,r)$. This excludes an actual multiplicity obstruction for
every padded-permanent closure, since $q\leq a=K_n$.

Indeed, $\lambda=\mu+d(n-b)e_1$, so O-012.3 gives
$K_n(\lambda)\geq a_\mu(d[b])=a_\lambda(d[n])$.
The reverse inequality is the ambient coordinate-ring quotient.

In particular, suppose $m\geq2$, $n\geq m^{25}$, and
$r\leq m^2$. Every shape satisfying $b\leq25$ in O-012.c
has full ambient determinant multiplicity. To verify the uniform
bound, use monotonicity in $b,r$:
\[
 B(b,r)\leq B(25,m^2)=2\binom{m^2+11}{12}-1.
\]
For $m=2$ the latter is $909<2^{25}$. For $m\geq3$,
\[
 2\binom{m^2+11}{12}
  \leq\frac{2(3m^2)^{12}}{12!}
  <m^{24}\leq m^{25},
\]
because $m^2+11\leq3m^2$ and $2\cdot3^{12}<12!$.
Also $n\geq m^{25}\geq25\geq b$.

For the explicit first higher-multiplicity example, take $d\geq6$,
$n\geq8$ and
\[
               \lambda=(d(n-2)-4,8,2^{d-2}).
\]
Here $k=2$, $t=6$, and $b=8$. Removing the column
$\det^2$ factor and using the binary weight count of F-013
(equivalently O-006.b) gives
\[
 a_\lambda(d[n])=p_6(d,n-2)-p_5(d,n-2)=11-7=4.
\]
Both rectangle bounds are inactive since $d,n-2\geq6$.
Therefore
\[
       n\geq\max\{8,B(8,d)\}\quad\Longrightarrow\quad
                         K_n(\lambda)=a_\lambda(d[n])=4. \tag{O-012.7}
\]
In the BIP range this holds for every $6\leq d\leq m^2$.
The statement settles all four determinant-side copies and
excludes $q>K_n$ there; it does not assert a value for the
permanent-side multiplicity between zero and four.

\paragraph{Geometric dictionary and scope.}
The determinant matrices define morphisms of affine coefficient
schemes over $\mathbb Z$. Multiplication by a chosen section
and contraction by its dual are maps of locally free symmetric
powers; the factorial diagonal calculation specifies their
characteristic-zero effect on all highest-weight copies.
The stabilization is a calculation in the coordinate algebra
of the matrix scheme: a determinant-line factor, followed by
division by a common first-row monomial, preserves the indicated
torus-weight and finite-group-invariant modules. The precise
background is the Cauchy identity, Macdonald I(4.3), p.~63,
and the functorial Schur decomposition I Appendix~A(5.3)--(5.6),
pp.~155--156, together with SGA~3 I4.7.2--4.7.3 and SGA~6
VI3.3. The related known padding statement is IP Proposition~2.6(b).
No claim is made that the factorial factors are invertible over
all of $\mathbb Z$. Integral polynomial identities and
characteristic-zero irreducible multiplicities remain separate.

The new gate is exact: a shape in this stabilized class can
only be an actual multiplicity obstruction if it fails
(O-012.5). In particular, at $n\geq m^{25}$ a surviving
maximal-row shape $\ell(\lambda)=d$ must have
$\lambda_d+\sum_{i=2}^d(\lambda_i-\lambda_d)>25$.
This is a full-multiplicity exclusion, including non-scalar
ambient multiplicities, rather than another occurrence test.
The later quadratic Chow bound M-010 excludes all maximal-row
shapes with $d\leq m^2$ throughout the BIP range as well.
Thus this particular BIP specialization is superseded by that
stronger degree result. O-012.a and O-012.b remain the inputs
to F-014 and F-015, which apply also at arbitrarily large outer
degree; the parameter-dependent statement (O-012.5) remains
valid in every stated range.
